\chapter{Introdução}
%\addcontentsline{toc}{chapter}{Introdução}

Hospitais são instituições fundamentais para a manutenção da saúde pública, enfrentando uma demanda crescente por seus serviços, impulsionada por fatores como o aumento e o envelhecimento populacional \cite{brandeau2004operations, schafer2023combining}. Contudo, essa demanda ocorre em um cenário de alta pressão por eficiência, devido aos custos ascendentes dos serviços de saúde e à limitação de recursos financeiros e físicos \cite{brandeau2004operations, schafer2023combining}. Nesse contexto, a gestão eficiente dos recursos hospitalares não é apenas uma meta operacional, mas uma necessidade estratégica para garantir a sustentabilidade do sistema e a qualidade do atendimento ao paciente \cite{brandeau2004operations, he2019systematic}.

No Brasil, onde o Sistema Único de Saúde (SUS) enfrenta desafios históricos de superlotação e subfinanciamento, a otimização de leitos torna-se ainda mais crucial. Um dos recursos mais críticos e, frequentemente, o gargalo das operações hospitalares, é a disponibilidade de leitos de internação \cite{brandeau2004operations, he2019systematic}. A gestão ineficiente de leitos gera impactos em cascata por todo o hospital, sendo apontada como uma das principais causas da superlotação nos serviços de emergência, um fenômeno conhecido como \textit{bed-block} ou ``bloqueio de leitos'' \cite{brandeau2004operations, he2019systematic}. Essa superlotação resulta em longos tempos de espera para pacientes admitidos, desvio de ambulâncias para outras unidades e, consequentemente, compromete a segurança e a qualidade do cuidado \cite{brandeau2004operations, he2019systematic}.

A complexidade do problema é acentuada pela natureza estocástica do ambiente hospitalar, marcada pela incerteza nas chegadas de pacientes, especialmente os de emergência, e pela variabilidade no tempo de permanência (LOS - \textit{Length of Stay}) de cada paciente \cite{ceschia2012modeling, he2019systematic, schafer2023combining}. Além disso, a alocação de pacientes a leitos transcende uma simples otimização de ocupação. Trata-se de um desafio multiobjetivo que deve conciliar interesses conflitantes de múltiplos \textit{stakeholders}, como pacientes, profissionais da saúde e a administração hospitalar \cite{schafer2023combining}. Este trabalho, contudo, trata a versão determinística do problema, na qual as datas de admissão e de alta de cada paciente são conhecidas.

Do ponto de vista do paciente, busca-se adequação clínica e conforto, como a minimização de transferências internas \cite{schafer2023combining}. Os profissionais da saúde e a administração têm ainda interesses próprios, como o equilíbrio da carga de trabalho das equipes e a utilização eficiente dos recursos \cite{ceschia2011local, demeester2010hybrid, schafer2023combining}. Entre esses critérios, este trabalho considera a adequação da especialidade clínica, as transferências, a separação de pacientes de gêneros distintos e a capacidade dos quartos.

Para superar os métodos tradicionais, diversas abordagens de Pesquisa Operacional, como programação matemática, simulação e modelos baseados em teoria das filas, têm sido propostas para enfrentar esse desafio \cite{brandeau2004operations, he2019systematic}. Métricas como taxa de ocupação, intervalo de substituição entre pacientes (\textit{turnaround time}) e tempo médio de permanência (\textit{length of stay}) são frequentemente utilizadas na literatura para avaliar a eficiência da alocação de leitos \cite{he2019systematic}.
No entanto, o problema de alocação de leitos (PBA - \textit{Patient-Bed Allocation} ou PAS - \textit{Patient Admission Scheduling}) é reconhecido por sua alta complexidade computacional, sendo frequentemente classificado como NP-difícil ou NP-completo \cite{ceschia2012modeling, demeester2010hybrid, schafer2023combining}. Isso significa que encontrar uma solução ótima por meio de métodos exatos (como a programação linear inteira) é, na maioria das vezes, inviável para hospitais de grande porte em tempo hábil \cite{ceschia2012modeling, demeester2010hybrid}.

Mesmo em sua versão determinística, em que as datas de admissão e de alta são conhecidas, o problema cresce rapidamente com o porte do hospital, e os métodos exatos deixam de produzir boas soluções em tempo compatível com o planejamento rotineiro do censo hospitalar \cite{ceschia2012modeling, demeester2010hybrid}. Essa lacuna justifica o desenvolvimento de algoritmos heurísticos e meta-heurísticos, capazes de fornecer soluções de boa qualidade em tempo compatível com esse planejamento, mesmo em instâncias de grande porte \cite{demeester2010hybrid, talbi2009metaheuristics}.

Diante do exposto, o objetivo central deste trabalho é o desenvolvimento de um algoritmo heurístico para o planejamento tático do censo hospitalar, isto é, a alocação de um conjunto conhecido de pacientes, com datas de admissão e de alta previamente definidas, aos quartos do hospital ao longo de um horizonte de planejamento. O problema é formulado como um modelo multiobjetivo de Programação Linear Inteira, sendo a função objetivo tratada como uma soma ponderada de penalidades. A admissão em tempo real, com chegadas de emergência e incerteza no tempo de permanência, está fora do escopo deste trabalho.

Para contextualizar os resultados obtidos pela heurística, também foi implementada uma abordagem exata, na qual o mesmo modelo é resolvido pelo \textit{solver} comercial Gurobi. Ambas as abordagens são comparadas quanto à qualidade da solução e ao esforço computacional em instâncias sintéticas de diferentes portes, de modo a identificar seu comportamento à medida que o tamanho do cenário simulado aumenta.

O modelo proposto, comum às abordagens heurística e exata, considera variáveis binárias que indicam em qual quarto cada paciente permanece em cada dia de sua internação, sendo as datas de admissão e de alta parte dos dados de entrada. Esse modelo tem como objetivo minimizar simultaneamente as alocações em quartos de especialidade inadequada, as transferências durante a internação e a ocupação concomitante de quartos por pacientes de gêneros distintos, respeitando a capacidade de leitos de cada quarto.

Por fim, este trabalho se propõe a criar uma ferramenta de apoio ao planejamento do censo que mantenha soluções de boa qualidade mesmo em instâncias de grande porte, típicas de grandes instituições hospitalares, nas quais os métodos exatos deixam de ser viáveis no tempo disponível para esse planejamento.